\chapter{Introdução}
\label{chap:intro}

Automodelos radiocontrolados de competicao possuem sistemas de propulsao de
elevada potencia especifica, porém operam tipicamente em malha aberta: o
operador envia um comando de aceleracao via radio e o controlador eletronico de
velocidade (ESC) aplica ao motor uma tensao correspondente, sem qualquer
realimentacao das variaveis de estado do sistema. No veiculo utilizado neste
projeto, um monster truck eletrico Kyosho Mad Force Kruiser 2.0 VE com motor
brushless de 2100 KV e ESC de 120 A, essa configuracao impede que estrategias
de controle dinamico de velocidade em malha fechada sejam aplicadas e deixa o
sistema desprotegido contra esforcos mecanicos e eletricos excessivos.

A aceleracao do veiculo influencia diretamente sua resposta transitoria: o
controle de velocidade define o valor final a ser atingido, enquanto o controle
de aceleracao define o perfil da curva, determinando a rapidez e a suavidade
com que o regime de operacao e alcancado. Esse tipo de controle relaciona-se
diretamente aos problemas classicos de seguimento de referencia (PSR) e de
rejeicao de perturbacoes (PRP), uma vez que atua sobre a dinamica do veiculo.
Alem disso, variacoes abruptas de esforco estao associadas a picos de corrente,
elevacao de temperatura e desgaste prematuro dos componentes eletromecanicos,
de modo que o monitoramento do consumo eletrico em tempo real tambem se
configura como mecanismo de protecao.

Este projeto propoe, desenvolve e valida a instrumentacao eletronica do
veiculo de forma a viabilizar o controle da aceleracao em malha fechada e o
monitoramento do consumo eletrico. Como etapa central, busca-se estabelecer
experimentalmente a curva caracteristica torque versus corrente do sistema de
propulsao completo, permitindo estimar indiretamente o torque entregue as
rodas a partir da corrente medida por um sensor de efeito Hall ACS758,
sem sensores mecanicos dedicados.

\section{Motivação e Justificativa}
\label{sec:motivacao}

O Departamento de Engenharia Eletronica (DELT) da UFMG oferece a disciplina
Veiculos Autonomos como optativa do curso de Engenharia de Controle e
Automacao, utilizando automodelos como plataformas experimentais. Em sua
configuracao original, o veiculo opera em malha aberta, o que limita seu uso
como ferramenta de ensino e pesquisa em controle. A instrumentacao proposta
transforma o automodelo em uma bancada movel para estudos de controle
dinamico, estimacao de variaveis mecanicas a partir de grandezas eletricas e
protecao de atuadores, consolidando-o como ferramenta experimental para a
disciplina e para grupos de pesquisa em robotica e sistemas inteligentes.

A estimacao de torque a partir da corrente e particularmente atraente por
dispensar torquimetros e dinamometros, equipamentos de elevado custo e de
difícil integracao ao conjunto compacto motor--transmissao do veiculo. A
caracterizacao da relacao torque--corrente para o sistema de propulsao
completo (motor, ESC, transmissao e perdas), em vez do motor isolado, ainda
é pouco documentada na literatura para plataformas de pequeno porte, o que
configura a contribuicao tecnica deste trabalho.

\section{Objetivos do Projeto}
\label{sec:objetivos}

Tendo em vista o exposto acima, este projeto tem por objetivos:

\begin{enumerate}[a)]
\item projetar e integrar a instrumentacao eletronica do automodelo
      (sensor de corrente, aquisicao embarcada e geracao de comandos),
      permitindo o monitoramento do consumo eletrico em tempo real;
\item desenvolver e executar um metodo experimental de baixo custo para
      obtencao da curva caracteristica torque versus corrente do sistema de
      propulsao completo, com tratamento de incertezas;
\item validar a curva obtida por concordancia com o modelo teorico do
      powertrain e quantificar suas caracteristicas estaticas
      (sensibilidade, linearidade, repetibilidade e histerese);
\item estabelecer as bases para o controle de aceleracao em malha fechada
      do veiculo, a ser implementado em trabalho futuro.
\end{enumerate}

\section{Local de Realização}
\label{sec:empresa}

As atividades foram desenvolvidas no Laboratorio de Sistemas Inteligentes
(LSI), localizado na Sala 2515 do Bloco 1 da Escola de Engenharia da
Universidade Federal de Minas Gerais, no campus Pampulha. O LSI e um
laboratorio pluridisciplinar e compartilhado, voltado ao desenvolvimento de
sistemas embarcados, automacao, robotica e controle para solucoes
inteligentes e autonomas.

\section{Estrutura da Monografia}
\label{sec:organizacao}

O trabalho esta dividido em cinco capitulos. Este capitulo apresentou a
introducao, a motivacao e os objetivos do projeto. O Capitulo 2 descreve o
sistema de propulsao do automodelo e a instrumentacao eletronica empregada,
incluindo a revisao das alternativas de sensores. O Capitulo 3 apresenta a
metodologia experimental para obtencao da curva caracteristica
torque--corrente. O Capitulo 4 apresenta os resultados experimentais e sua
analise, incluindo o tratamento de incertezas. O Capitulo 5 conclui o trabalho
e propoe continuacoes.

\clearpage
